% 
% ###################################################################
%  Copyright (c) 2025, AUTHORNAME 
%  Editors: A.D. Rollett & M. De Graef
%  All rights reserved.
%
% Licensed under the Creative Commons CC BY-NC-SA 4.0 License, 
% hereafter referred to as the "License"; you may not use this 
% document except in compliance with the License. You may obtain 
% a copy of the License at 
%     https://creativecommons.org/licenses/by-nc-sa/4.0/legalcode 
% Unless required by applicable law or agreed to in writing, all 
% material distributed under the License is distributed on an 
% "AS IS" BASIS, WITHOUT WARRANTIES OR CONDITIONS OF ANY KIND, 
% either express or implied. See the License for the specific 
% language governing permissions and limitations under the License.
% ###################################################################

% ###################################################################
% ###################################################################
% ###################################################################
% The following lines are to be uncommented or edited as needed 

%     uncomment this line only if this chapter is a core/foundational chapter;
%     for a core chapter a "Foundational" label will appear on the top left above the chapter title.
\corechapter{Yes}

%     this will appear in a secondary header below the chapter title.
\OCchapterauthor{Marc De Graef, Carnegie Mellon University}

% All figures are stored in the src/ChapterTitle/eps folder and must be of the *.pdf, *.pdf, or *.png type. 

% this is the 6-character (all uppercase) chapter descriptor used in labels and refs...
\renewcommand{\chabbr}{FOURIE}

%     replace \noheaderimage by the chapter header image file name (without .pdf extension).
%     pdf files are also acceptable.
%     Chapter header images must be 2480 x 1240 pixels with 300dpi, RGB format.
\chapterimage{\noheaderimage}

% start the chapter and define the chapter label (outside of the \chapter command!)
\chapter{Fourier Transforms and Series}\label{\chabbr:ch:FourierIntro} % should be the same as the filename (without extension)

% add the author information so that it will appear in the Author List at the start of the document
% repeat this line for each author; the first entry must match the chapter label above.
\writeauthor{\chabbr:ch:FourierIntro}{Fourier Transforms and Series}{De Graef}{Marc}{Materials Science and Engineering}{Carnegie Mellon University}{mdg@andrew.cmu.edu}{https://www.mse.engineering.cmu.edu/directory/bios/degraef-marc.html}

% Each chapter begins with Learning Objectives; the list of objectives should have links to sections/subsections using their OC labels.
% Each Learning Objective should be an active statement (i.e., contain a verb).  The \\ command can be used to force an item into the 
% second column if LaTeX breaks the line at an awkward location.  The label color (OCBurntOrange) should not be changed.
% Within the itemize environment, hyphenation is temporarily turned off.
\begin{learningobjectives}
% This itemized list should be formatted as follows and must be enclosed by { } 
%{\begin{itemize}
%    \item[{\color{OCBurntOrange}\ref{\chabbr:sec:complex}:}] Learn a formal description of
%\end{itemize}}
\end{learningobjectives}
% ###################################################################
% ###################################################################
% ###################################################################


%%%%%%%%%%%%%%%%%%%%%%%%%
%%%%%%%%%%%%%%%%%%%%%%%%%
% This where the actual chapter content starts   %
%%%%%%%%%%%%%%%%%%%%%%%%%
%%%%%%%%%%%%%%%%%%%%%%%%%

\section{Introduction}\label{\chabbr:sec:introduction}


\subsection{}\label{\chabbr:ssec:}



\section{Special Functions}\label{\chabbr:sec:SpecFunc}



\section{Fourier Series}\label{\chabbr:sec:FourierSeries}


\section{Fourier Transforms}\label{\chabbr:sec:FourierTransforms}


\section{Convolutions and Related Operations}\label{\chabbr:sec:Convolutions}












